\documentclass[10pt,a4paper]{amsart}
\usepackage[margin=19mm]{geometry}
\usepackage{amsmath,amssymb,mathtools}
\usepackage[scr=rsfs,cal=euler]{mathalfa}
\usepackage{xfrac}
\usepackage[unicode,hidelinks,bookmarks=false]{hyperref}
\newcommand{\ie}{i.e.\ }
\newcommand{\cf}{cf.\ }
\newcommand{\resp}{respectively\ }
\newcommand{\Subsection}{\subsection}
\providecommand{\nobreakdash}{\nobreak\hbox{-}}
\setlength{\emergencystretch}{2em}
\multlinegap0pt
\usepackage[shortlabels]{enumitem}
\usepackage{mathtools}
\usepackage{csquotes}
\usepackage{amssymb}
\usepackage{xcolor}
\usepackage{tikz}
\newtheorem{theorem}{Theorem}
\title{Thin set theorem for arbitrarily many colors implies bounding}
\author{Yamato Miyata, Keita Yokoyama}
\date{Source: arXiv:2608.25339v1 (CC BY 4.0)}
\begin{document}
\maketitle

\noindent\emph{Source and licence.} Thin set theorem for arbitrarily many colors implies bounding. Version arXiv:2608.25339v1 (CC BY 4.0), distributed by arXiv under the Creative Commons Attribution 4.0 International licence, http://creativecommons.org/licenses/by/4.0/, as recorded in the arXiv licence metadata for this exact version and shown on the abstract page https://arxiv.org/abs/2608.25339v1. Full licence notice: \texttt{licenses/arxiv-2608.25339-CC-BY-4.0.txt}.\par\smallskip

\noindent\textit{Editorial scope.} Counted unit: the article's theorem thm:equivalence (source0.tex lines 144--147). The article's complete original proof of it is delivered with it (source0.tex lines 150--262, one \texttt{proof} environment of 113 lines). No mathematics is written, repaired or re-derived: the statement and the proof are the article's own lines, in the article's own notation, and no retained line is substituted or re-flowed.  No further source-local context is needed: every cross-reference of every retained line resolves inside the two delivered spans.  Nothing was omitted from any retained span, so the package carries no omission record.  No retained line contains a citation, so the wrapper carries no bibliography and no bibliography entry.  No retained line contains a figure environment, an image-inclusion command or drawing code, so no image is delivered and the package declares no asset dependency.  Editorial formatting only, in addition: an AMS \texttt{amsart} preamble that restates the article's own declarations and macros used by the retained lines, namely the environments \texttt{theorem} on separate counters, together with the standard packages those lines need.  The retained environments print in this document as Theorem 1 in order of appearance, whereas the article prints the same items with the numbering of its own full document.  The complete native source of the article as distributed in the arXiv e-print for this exact version is bundled unchanged as \texttt{sources/208568/source0.tex}.
\begin{theorem}\label{thm:equivalence}
For any natural numbers $n,\ell \ge 1$,
\[\mathsf{RCA}_0\vdash\mathsf{B}\Sigma_{n+1}^{0}\leftrightarrow \Sigma_{n}^{0}\text{-}\mathsf{RT}_{<\infty}^{1-}\leftrightarrow\Sigma_{n}^{0}\text{-}\mathsf{RT}_{<\infty,\ell}^{1-}.\]
\end{theorem}

\begin{proof}[Proof of Theorem \ref{thm:equivalence}]
Here, recall that for any natural numbers $n \ge 1$, $\mathsf{B}\Pi_{n}^{0}$ is equivalent to $\mathsf{B}\Sigma_{n+1}^{0}$ over $\mathsf{RCA}_0$.
Then, by considering inductively, it suffices to show that for any natural numbers $n,\ell \ge 1$,
\[\mathsf{RCA}_0+\mathsf{B}\Sigma_{n}^{0}\vdash\mathsf{B}\Pi_{n}^{0}\leftrightarrow\Sigma_{n}^{0}\text{-}\mathsf{RT}_{<\infty}^{1-}\leftrightarrow\Sigma_{n}^{0}\text{-}\mathsf{RT}_{<\infty,\ell}^{1-}.
\]
The proof is carried out in two steps.


\begin{enumerate}[label=(\roman*),topsep=0.5em,itemsep=0.5em]
\item First, we prove the equivalence \[\mathsf{RCA}_0+\mathsf{B}\Sigma_{n}^{0}\vdash\mathsf{B}\Pi_{n}^{0}\leftrightarrow\Sigma_{n}^{0}\text{-}\mathsf{RT}_{<\infty}^{1-}.\]


For the backward direction $(\leftarrow)$, we proceed by contradiction. Assume $\Sigma_{n}^{0}\text{-}\mathsf{RT}_{<\infty}^{1-}\land\neg\mathsf{B}\Pi_{n}^{0}$. Then, by $\neg\mathsf{B}\Pi_{n}^{0}$, there exists a $\Pi_{n}^{0}$ formula $\forall z \, \theta(x,y,z)$ (where $\theta$ is a $\Sigma_{n-1}^{0}$ formula) and a witness $w$ satisfying the following two conditions:
\begin{equation}
\forall x < w \, \exists y \, \forall z \, \theta(x,y,z) \label{lem:equivalence-1}
\end{equation}


\begin{equation} \forall v \, \exists x < w \, \forall y < v \, \exists z \, \neg\theta(x,y,z) \label{lem:equivalence-2}
\end{equation}
Let $v$ be arbitrary. By \eqref{lem:equivalence-2}, there exists an $x_0<w$ that satisfies the following:
\[
\forall y < v \, \exists z \, \neg\theta(x_0,y,z)
\]
Over $\mathsf{RCA}_0$, $\neg\theta$ is equivalent to a $\Pi_{n-1}^{0}$ formula. Thus, by $\mathsf{B}\Sigma_{n}^{0}$, the following bounds exist:
\[
\exists m \, \forall y < v \, \exists z < m \, \neg\theta(x_0,y,z)
\]
Thus, we have:
\begin{equation} \exists x < w \, \exists m \, \forall y < v \, \exists z < m \, \neg\theta(x,y,z) \label{lem:equivalence-3}
\end{equation}


We can then define a function $c(v)$ by:


\begin{equation} c(v)=\min\{ \langle x,m \rangle \mid x < w \land \forall y < v \, \exists z < m \, \neg\theta(x,y,z) \}.
\end{equation}
By \eqref{lem:equivalence-3} and the least number principle for $\Pi_{n-1}^{0}$ formulas ($\mathsf{L}\Pi_{n-1}^{0}$), which follows from $\mathsf{B}\Sigma_{n}^{0}$, $c$ is a total function.


Moreover, since $c$ is defined via bounded quantification and logical connectives applied to $\Sigma_{n-1}^{0}$ and $\Pi_{n-1}^{0}$ formulas, the induction/bounding hypothesis $\mathsf{B}\Pi_{n-1}^{0}$ ensures that $c$ can be expressed by a $\Sigma_{n}^{0}$ formula. Therefore, $c$ is a total $\Sigma_{n}^{0}$-function.


Next, we define $d\colon\mathbb{N}\to\mathbb{N}$ by
\[
d(v):=(c(v))_0.
\]
Since $d$ is the composition of $c$ and a total computable function, $d$ is also a total $\Sigma_{n}^{0}$-function. Furthermore, by the definition of $c$, $(c(v))_0<w$ holds, meaning that $d$ is a total function with a bounded range. Thus, we can apply $\Sigma_{n}^{0}\text{-}\mathsf{RT}_{<\infty}^{1-}$ to $d$, which gives:
\begin{equation}   \exists x < w \, \forall a \, \exists b > a \, \forall y < b \, \exists z \, \neg\theta(x,y,z)
\end{equation}
Then,
\begin{equation}
\exists x < w \, \forall y \, \exists z \, \neg\theta(x,y,z)
\end{equation}
However, this contradicts \eqref{lem:equivalence-1}. Therefore, by contradiction, the backward direction holds.


Next, we show the forward direction $(\rightarrow)$. Assume $\mathsf{B}\Sigma_{n+1}^{0}$. Fix $k$ and let $\varphi$ be an arbitrary total $\Sigma_{n}^{0}$-function $\varphi\colon\mathbb{N}\to\{0,1,\dots,k-1\}$. We will show by contradiction that there exists a color that appears infinitely often. Suppose for contradiction that this does not hold, i.e.,
\[
\forall i < k \, \exists x \, \forall z \, (z > x \rightarrow \neg\varphi(z,i)).
\]
Since the subformula $\forall z \, (z > x \rightarrow \neg\varphi(z,i))$ is $\Pi_{n}^{0}$, $\mathsf{B}\Sigma_{n+1}^{0}$ implies:
\[
\exists v \, \forall i < k \, \exists x < v \, \forall z \, (z > x \rightarrow \neg\varphi(z,i)).
\]
In particular, this implies:
\[
\exists v \, \forall i < k \, \forall z > v-1 \, \neg\varphi(z,i).
\]
Then,
\[
\exists v \, \forall i < k \, \neg\varphi(v,i).
\]
However, this contradicts our assumption that the range of $\Sigma_{n}^{0}$-function $\varphi$ is contained in $\{0,1,\dots,k-1\}$. Thus, by contradiction, $\Sigma_{n}^{0}\text{-}\mathsf{RT}_{<\infty}^{1-}$ holds, which completes the proof of $\mathsf{RCA}_0+\mathsf{B}\Sigma_{n}^{0}\vdash\mathsf{B}\Pi_{n}^{0}\leftrightarrow\Sigma_{n}^{0}\text{-}\mathsf{RT}_{<\infty}^{1-}$.


\item Finally, we show that for all natural numbers $\ell\ge 1$,
\[ \mathsf{RCA}_0+\mathsf{B}\Sigma_{n}^{0}\vdash\Sigma_{n}^{0}\text{-}\mathsf{RT}_{<\infty,\ell}^{1-}\leftrightarrow\Sigma_{n}^{0}\text{-}\mathsf{RT}_{<\infty}^{1-}.
\]
Since the direction $(\leftarrow)$ is trivial, we focus on $(\rightarrow)$. Suppose $\Sigma_{n}^{0}\text{-}\mathsf{RT}_{<\infty,\ell}^{1-}$ holds. For any $k$ and total $\Sigma_{n}^{0}$-function $\varphi\colon\mathbb{N}\to\{0,1,\dots,k-1\}$, we have:


\begin{equation}
\exists s \subseteq \{0,1,\dots,k-1\} \, (|s|\le\ell \land \forall a \, \exists b \, [a < b \land \exists i \in s \, \varphi(b,i)]).
\end{equation}


We assume $\Sigma_{n}^{0}\text{-}\mathsf{RT}_{<\infty}^{1-}$ fails. Then, there exist $k$ and a total $\Sigma_{n}^{0}$-function $\varphi\colon\mathbb{N}\to\{0,1,\dots,k-1\}$ such that:


\begin{equation}
\forall i < k \, \exists a \, \forall b \, (a < b \rightarrow \neg\varphi(b,i)). \label{lem:equivalence-4}
\end{equation}


We fix those $\varphi$ and $\ell$ and let $s \subseteq \{0,1,\dots,k-1\}$ with $|s|\le\ell$ be the set obtained by applying $\Sigma_{n}^{0}\text{-}\mathsf{RT}_{<\infty,\ell}^{1-}$ for $\varphi$.
Then, for each $i \in s$, let $a_i$ be a witness satisfying the condition of \eqref{lem:equivalence-4}. Now, we define $m:=\max\{a_i \mid i\in s\}$. (Since the size of $s$ is at most $\ell$, the existence of $m$ can be demonstrated within $\mathsf{RCA}_0$.)
Then the following holds:


\begin{equation}
\forall m' > m \, \forall i \in s \, \neg\varphi(m',i). \label{lem:equivalence-5}
\end{equation}


On the other hand, by the assumption of $\Sigma_{n}^{0}\text{-}\mathsf{RT}_{<\infty,\ell}^{1-}$, there must exist a value such that:
\[
\exists m' > m \, \exists i \in s \, \varphi(m',i)
\]
which contradicts \eqref{lem:equivalence-5}. Thus, by contradiction, we obtain $(\rightarrow)$.
\end{enumerate}
\end{proof}

\end{document}
